\section{Weekly Summary}

\begin{tcolorbox}[colback=yellow!10,colframe=black!50,title=AI-Generated Content]
\noindent The summary below was written by Claude (Anthropic) from that week's regression outputs and series descriptions. It has not been reviewed or fact-checked by a human, and should be read as an automated, informal recap -- not analysis.
\end{tcolorbox}

Welcome back to the weekly ritual where we throw two unrelated economic time series into a blender, press the button, and act surprised at whatever sludge comes out. This week's lineup featured light-truck seasonal adjustment factors, Latin American fish prices, the minimum wage in our nation's capital, Midwestern gasoline, Japanese producer prices, and some discontinued data about how much cash the upper-middle-class keeps lying around. A real murderer's row of economic indicators that have absolutely no business being in the same room together. Naturally, we introduced them anyway.

Monday kicked us off with a resounding shrug. The seasonal factors for light truck production versus nonferrous metal export prices produced an R-squared of 0.00003, which is statistics-speak for "these two have never met and never will." The p-value was 0.93, a number so comfortably insignificant it could be framed and hung in a museum of things that don't matter. No trend to blame, no relationship to find, nothing collapsing when detrended because there was never anything standing in the first place. A beautiful, pure null result. We salute it. Tuesday wasn't much better: Latin American agriculture import prices and the cash holdings of the 50th-to-90th wealth percentile looked impressively linked in the raw regression (p-value 0.0006, dramatic gasp), but the moment we stripped out the shared time trend, the whole thing evaporated into a p-value of 0.86. Textbook trend mirage. Two series drifting along in the same general direction and fooling exactly nobody who bothered to look twice.

Then Thursday arrived with the week's headline act: the District of Columbia minimum wage versus the Peoria, Illinois house price index, posting a raw R-squared of 0.77 and a p-value with so many zeros it needed scientific notation to describe its enthusiasm. For one glorious moment you might think DC's wage floor is personally propping up the Peoria housing market. And then, as always, we detrend -- and the relationship face-plants to a p-value of 0.91. Turns out both numbers just go up over time, like everything else in an economy, and that shared upward drift was doing all the heavy lifting. Another spectacular correlation reduced to rubble by the cruel arithmetic of a linear trend.

The one that actually survived came Friday: Midwest gasoline prices versus Japanese industrial producer prices. The raw regression was significant (no shock there), but it stayed significant even after detrending, with a p-value around 0.025. Is this the rare, interesting case? Possibly! Energy prices feeding into producer prices across the globe is at least a story a person could tell with a straight face at a dinner party. Or it's noise wearing a convincing costume. As always, the mandatory reminder: these are randomly paired series, none of this implies causation, and drawing real conclusions from this project would be like reading tea leaves, except the tea is FRED and the leaves are a p-value. Spurious is the house style. See you next week.
